\documentclass[10pt]{amsart}

\usepackage[T1]{fontenc}
\usepackage{lmodern}
\usepackage{microtype}
\usepackage[a4paper,margin=0.82in]{geometry}
\usepackage{amsmath,amssymb,amsthm}
\usepackage[hidelinks]{hyperref}

\newtheorem{theorem}{Theorem}
\newtheorem{corollary}[theorem]{Corollary}
\theoremstyle{remark}
\newtheorem{remark}[theorem]{Remark}

\newcommand{\ldir}{\ell^{\mathrm{dir}}}
\newcommand{\kdir}{k^{\mathrm{dir}}}

\title{A Short Bound for Directed Local Connectivity}
\author{}
\date{}

\begin{document}

\begin{abstract}
A theorem of Huang and Lyu on repeated directed two-paths immediately
improves the known error term for digraphs of bounded local connectivity.
For every fixed $m\ge3$, both the arc- and vertex-disjoint extremal numbers
are at most $n^2/4+(m-1)n+O_m(1)$.
\end{abstract}

\maketitle

Let $P_{s,2}$ denote the union of $s$ directed paths of length two with
common initial and terminal vertices.  Huang and Lyu proved the following
sharp estimate for sufficiently large order.

\begin{theorem}[Huang--Lyu \cite{HL}]\label{thm:HL}
For $t\ge2$ and
\[
 n\ge N(t):=\max\left\{t^3+4t^2+3t+4,
 \frac{17}{2}t^2+30t+27\right\},
\]
every $P_{t+1,2}$-free simple digraph has at most
\begin{equation}\label{eq:HL}
 G(n,t):=
 \left\lceil\frac{n+t}{2}\right\rceil
 \left\lfloor\frac{n-t}{2}\right\rfloor+tn+1
\end{equation}
arcs.
\end{theorem}

For a simple digraph $D$, let $\lambda_D(u,v)$ be the maximum number of
arc-disjoint directed $u$--$v$ paths and let $\kappa_D(u,v)$ be the maximum
number of internally vertex-disjoint such paths.  Write $\ldir_m(n)$ and
$\kdir_m(n)$ for the corresponding extremal arc counts under the ceiling
$m-1$.

\begin{corollary}\label{cor:main}
Let $m\ge3$ and $n\ge N(m-1)$.  Then
\begin{align*}
 \ldir_m(n)\le\kdir_m(n)
 &\le G(n,m-1)\\
 &=\left\lceil\frac{n+m-1}{2}\right\rceil
   \left\lfloor\frac{n-m+1}{2}\right\rfloor
   +(m-1)n+1.
\end{align*}
Consequently, for each fixed $m\ge3$,
\[
 \ldir_m(n),\ \kdir_m(n)
 \le \frac{n^2}{4}+(m-1)n+O_m(1).
\]
\end{corollary}

\begin{proof}
If $D$ contained $m$ directed two-paths $u\to x_i\to v$ with common
endpoints, then their distinct internal vertices would make the paths both
arc-disjoint and internally vertex-disjoint.  Thus either connectivity
ceiling makes $D$ $P_{m,2}$-free, and Theorem~\ref{thm:HL} applies with
$t=m-1$.  Finally $\kappa_D(u,v)\le\lambda_D(u,v)$ gives
$\ldir_m(n)\le\kdir_m(n)$, while
\[
 \left\lceil\frac{n+t}{2}\right\rceil
 \left\lfloor\frac{n-t}{2}\right\rfloor
 \le \frac{n^2-t^2}{4}
\]
gives the asymptotic form.
\end{proof}

\begin{remark}
The augmented bipartite construction gives
\[
 \ldir_m(n)\ge
 \left\lfloor\frac{(n+m-2)^2}{4}\right\rfloor
 =\frac{n^2}{4}+\frac{m-2}{2}n+O_m(1).
\]
Hence Corollary~\ref{cor:main} reduces the previous upper coefficient
$4(m-1)$ to $m-1$, but it does not determine the conjectured coefficient
$(m-2)/2$: the two-path reduction deliberately discards all information
about direct arcs and longer routes.
\end{remark}

\begin{thebibliography}{1}
\bibitem{HL}
Z.~Huang and Z.~Lyu,
\emph{Extremal digraphs containing at most $t$ paths of length $2$ with the
same endpoints}, Discrete Appl. Math. \textbf{388} (2026), 1--10.
\end{thebibliography}

\end{document}
